\documentclass[12pt,a4paper]{article}
\usepackage[a4paper,margin=1in]{geometry}
\usepackage{graphicx}
\usepackage{hyperref}
\usepackage{titlesec}
\usepackage{enumitem}
\usepackage{longtable}
\usepackage{booktabs}
\usepackage{fancyhdr}

%----------------------------
% Title and Formatting
%----------------------------
\titleformat{\section}{\large\bfseries}{\thesection}{1em}{}
\titleformat{\subsection}{\normalsize\bfseries}{\thesubsection}{1em}{}
\setlist{nosep}

\pagestyle{fancy}
\fancyhf{}
\rhead{Technical Design Document}
\lhead{\leftmark}
\cfoot{\thepage}

\title{\textbf{Technical Design Document}\\[0.5em]\large Nimbler Sentences With Reinforcement Learning}

\author{Nimble @ WashU Robotics \\ \texttt{washurobotics@gmail.com}}

\author{
    \\
    Tyler Dinh, Lauren Gayle, Anling Zhu, Tenzin Dukdak, \\
    Billy Zhao, Noah Wolk, Mark Seidel, Horlasy Demakpor \\[0.3em]
    Software Team \\[0.2em]
    \texttt{\{dinh.t, gayle.l, anling, dtenzin, billy.z,} \\
    \texttt{n.j.wolk, m.m.seidel, horlasy.d\}@wustl.edu} \\\\
    Elizabeth Joseph, Alexis Luna, Yuxuan Shen, \\
    Kyle Okui, Gabe Minton \\[0.3em]
    Integration Team \\[0.2em]
    \texttt{\{elizabeth.j, alexis.l, notsure, k.okui, g.m.minton\}@wustl.edu}\\\\
}
\date{\today}

%----------------------------
\begin{document}
\maketitle
\newpage
\tableofcontents
\newpage

%----------------------------
\section*{Action Items}
% This is to be removed on completion of this document.
\subsection*{Software}

\begin{enumerate}

  \item \textbf{LeRobot Dataset Collection} \\
  Collect a \texttt{LeRobotDataset} for eventual VLA fine-tuning. This is the highest priority task.
  \begin{itemize}
    \item Data points will be collected through Isaac Lab or the Trossen Data Collection Framework.
    \item Prior to collection, the Isaac Lab virtual environment / physical environment must be configured to match the physical environment as closely as possible, including arm attachments (Logitech camera), and environmental objects (whiteboard, marker holder, etc.).
    \item The raw HDF5 dataset will first be released publicly on HuggingFace, followed by the converted \texttt{LeRobotDataset}.
  \end{itemize}

  \item \textbf{VLA Fine-Tuning} \\
  Fine-tune NORA 1.5 for handwriting tasks. Zero-shot inference of the NORA 1.5 VLA has been observed to not perform well in this domain, so we will fine-tune it.
  \begin{itemize}
    \item NORA 1.5's pretraining data likely contains no robotic writing examples, making fine-tuning essential for this specific scenario of writing.
    \item Fine-tuning will use the HuggingFace \texttt{LeRobotDataset} via the fine-tuning script provided in the NORA 1.5 GitHub repository.
    \item The fine-tuned VLA will be released publicly on HuggingFace. This work has potential for academic writing, like a technical white paper or report given the project's structure. The fine-tuned model could also be shared on social media to gauge interest from researchers and practitioners. A faculty collaborator may be worth pursuing in the future after we have completed this project.
  \end{itemize}

\end{enumerate}

\subsection*{Integration}

\begin{enumerate}

  \item \textbf{Arm and Camera Wiring} \\
  Physical wiring of the arm and camera into the compute environment must be established and discussed with the team. Since the Roboserver will be running VLA inference and all associated scripts, a reliable communication channel between the arm, camera, and Roboserver is required.

  \item \textbf{Physical Environment Setup} \\
  The physical environment of the arm must match the Isaac Lab virtual environment used for data collection and fine-tuning as closely as possible. Mismatches between virtual and physical environments are a primary source of performance degradation.

  \item \textbf{Linear Actuator Integration} (\textit{in progress}) \\
  A linear actuator enables the arm to traverse the writing surface horizontally, extending its effective writing range. Without this, the robot cannot produce arbitrarily long text on a surface larger than its static range of motion.

\end{enumerate}
\newpage
%----------------------------

%----------------------------
\section{Overview}
%Describe the purpose of this document and its intended audience.  
%Mention any prerequisite knowledge required before reading.  
%Specify whether it targets a technical or non-technical audience.
The purpose of this document is to present the technical design of a system that enables a robotic arm to learn and perform handwriting tasks in a physical environment for basic sentences using reinforcement learning (RL). \\\\
This document is intended for a technical audience with a background in computer science, robotics, or machine learning. Non-technical readers may find high-level summaries useful but should refer to supplementary documentation at the end of this document for background on the underlying mechanisms employed.
\newpage
%----------------------------
\section{Problem}

Current robot arms used for writing tasks operate with \textbf{static, pre-programmed motions}. Each specific writing task requires its own custom motion program, which has several critical limitations:

\subsection{The Core Problem}

If a robot needs to write something new, engineers must:
\begin{itemize}
    \item Manually program the exact trajectory for each letter or word.
    \item Hard-code the sequence of movements.
    \item Test, refine, and repeat this entire process for every new writing task.
\end{itemize}
An autonomous robot that generalizes its writing capabilities would:
\begin{itemize}
    \item Scale efficiently: Write any arbitrary text without reprogramming.
    \item Adapt to variations: Handle different sizes and styles of text.
    \item Learn from examples: Improve through demonstration rather than explicit programming.
\end{itemize}

\subsection{The Core Goal}

The goal is to develop a robot writing system that can:
\begin{itemize}
    \item Accept arbitrary text input (like ``Hello World'' or even more arbitrary, like ''Write me a haiku.'')
    \item Dynamically generate appropriate writing motions.
    \item Execute smooth, legible handwriting without pre-programmed trajectories for each specific string.
\end{itemize}
\newpage
%----------------------------
\section{Requirements}
TODO! Describe all requirements imposed by the problem.  
Write from the end-user’s perspective. Include user stories or use cases.

\subsection*{Example Use Cases}
\begin{itemize}
  \item As a x, I want to y.
  \item As a x, I want y.
\end{itemize}
%----------------------------
\subsection{Out of Scope}
TODO! List what is explicitly out of scope to avoid misunderstandings.

%----------------------------
\subsection{Success Criteria}
TODO! Describe how success will be measured once the solution is in production.  
Include measurable metrics such as performance, scalability, or cost efficiency.

\subsection*{Example Metrics}
\begin{itemize}
  \item Time-to-market reduced by 90\%.
  \item Supports 10,000 users with 100ms P90 latency.
  \item It can write sentences in English legibly with X\% accuracy.
\end{itemize}
\newpage
%----------------------------
\section{Architecture}
Describe the architecture of the proposed solution using text, diagrams, and bullet points.  
Explain why this design was chosen.

25-26 Architecture

VLA: NORA 1.5
Data Collection Framework: LeRobot
Arm: WidowX AI Trossen Base
CV: OpenCV + Logitech
Additional VLM/LLM for Data Preprocessing: Local Qwen3-4B instance,
Programming Languages: Mainly Python, JSON for data.
Compute: RoboServer

\subsection{VLA Dataset and Training}

During operation, the VLA will be asked by the user to write a sentence or word in the format of \texttt{Write `THING\_TO\_WRITE'}.

\begin{description}
  \item[Q: How do we ensure the VLA can understand all text inputs?]
  We need to collect lots of diverse data.

  \item[Q: How do we collect diverse data for the text of a sentence?]
  We need to learn how to write characters in that language so that text can be broken down into atomic units.
\end{description}

For collecting diverse data points for writing English text, we define a control script to write characters given their stroke-by-stroke physical definitions. Each data point for a character contains:

\begin{enumerate}
  \item The writing instruction: \texttt{Write `character'}
  \item The vision: camera input captured before writing that stroke
  \item The stroke movement: the arm trajectory required to draw the next stroke of the desired character
\end{enumerate}

\begin{description}
  \item[Q: Why not just train on words?]
  VLAs have sparse rewards, only receiving a reward signal at the end of the sequence. Writing a word takes many timesteps. More fundamentally, the VLA outputs a single 6-DOF action vector per inference pass---there is no single action that writes a word. Writing even a short word requires hundreds of sequential arm poses, plus pen-lifts within characters (e.g., the dot in `i'), between characters, and between words. Training at the character level avoids all of these issues.
\end{description}

\subsection{Fine-Tuned VLA Inference}

A sentence is decomposed hierarchically:
\begin{itemize}
  \item A \textbf{sentence} decomposes into words, separated by scripted spaces.
  \item A \textbf{word} decomposes into characters, separated by scripted smaller spaces.
  \item A \textbf{character} decomposes into strokes. A \textbf{stroke} is the atomic unit of VLA inference.
\end{itemize}

For example, `A' has 3 strokes and `I' has 1 stroke. The VLA performs one inference pass per stroke, and the stroke count per character is physically defined in code, so the system always knows exactly how many inference passes to run. Each inference pass receives:
\begin{enumerate}
  \item The writing instruction for the \emph{character}: \texttt{Write `character'}
  \item The vision: camera input captured before executing that stroke
  \item Output: the stroke movement vector used to actuate the arm
\end{enumerate}

Pen-lifts and repositioning between characters and words are \textbf{scripted}, not learned, so they do not require VLA inference.

\subsubsection*{Control Flow}

\begin{verbatim}
"Write 'hello'"
-> parse -> ['h', 'e', 'l', 'l', 'o']
-> for each character:
   -> look up stroke count in physically defined constant
   -> run VLA inference stroke_count times
   -> scripted pen-lift + reposition to next character start position
-> done

"Write 'hello world'"
-> parse -> [['h','e','l','l','o'], ['w','o','r','l','d']]
-> for each word:
   -> for each character:
      -> look up stroke count in physically defined constant
      -> run VLA inference stroke_count times
      -> scripted pen-lift + reposition to next character start position
   -> scripted larger pen-lift and reposition for word gap
-> done
\end{verbatim}

\subsection{Extension of Writing Range with a Linear Actuator}

Writing text is difficult, as the phrase `hello world this is a cool sentence` with the WidowX arm can only be written in the range of action given by the robot and its camera. A linear actuator provides a `sliding window' over the writing surface, enabling arbitrarily long text output.

\subsubsection*{Text Planning}

Before any motion occurs, the full spatial layout of the text should be pre-computed:
\begin{itemize}
  \item Determine character width and inter-character spacing in physical units (cm).
  \item Map the full string to an absolute coordinate sequence along the writing axis.
  \item Determine the total horizontal distance required.
\end{itemize}

\subsubsection*{Text Chunking by Reachability}

Given the current position of the arm, a \textbf{range of action} is defined, which is the zone the writing tool can write in at the current position. The input text to the start of the program is greedily chunked, meaning that characters are grouped until the next character would fall outside the window. Chunk boundaries are placed \emph{between characters}, but should never be chunked mid-character.

\subsubsection*{Slide by Inference on a Chunk}

For each chunk, the VLA (NORA 1.5) executes the writing motions normally. When the chunk is complete, the actuator must activate, moving X distance based on where the next chunk is. Actuator movement is intentionally seperate from VLA inference, as it is purely physical and does not need to be learned. After sliding, a brief re-anchoring step determines where the last written character ended to determine where the arm should resume movement. The VLA is then re-prompted with the next chunk of text.

\subsection{High-Level Overview (HLD)}
List all logical system components:
\begin{itemize}
  \item comp1
  \item comp2
  \item comp3
  \item comp4
  \item comp5
\end{itemize}

Include diagrams as needed:
\begin{center}
%\includegraphics[width=0.8\textwidth]{architecture-diagram.png}
\end{center}

\subsection{API Design}
List all APIs used for interaction between users or services.  
Include HTTP methods, payloads, versions, and example requests/responses.  
Discuss future evolution of the APIs.

\subsection{Data Storage and Model}
Describe the data model and database choice.  
Estimate data volume, forecast growth, and justify scalability.  
Discuss data pipelines, ingestion, and pre-processing if applicable.

\subsection{Application / Component Level Design (LLD)}
Provide detailed design of each system component, including data flow and control flow diagrams.

USE DRAW.IO for this (can install on your computer)

\newpage
%----------------------------
\section{Dependencies}
List external and internal dependencies.  
Document assumptions and risks associated with these dependencies.
%----------------------------
\subsection{Design Alternatives Considered}
Discuss alternative designs evaluated.  
Provide a comparison table highlighting trade-offs.

\begin{longtable}{@{}p{3cm}p{5cm}p{5cm}@{}}
\toprule
\textbf{Option} & \textbf{Pros} & \textbf{Cons} \\ \midrule
Option A & Scalable, cost-efficient & Complex to maintain \\
Option B & Easy to implement & Limited scalability \\ \bottomrule
\end{longtable}

\newpage
%----------------------------
\section{Cost Analysis}
Estimate infrastructure and operational costs.  
Plan for future growth and scalability.

\newpage
%----------------------------
\section{Failures and Risks}
Discuss potential system failures such as:
\begin{itemize}
  \item Dependency failures
  \item Traffic overflow
  \item Performance degradation
  \item Logic bugs
\end{itemize}

Identify risks such as external dependencies or resource limitations.  
Include potential mitigation strategies.

\newpage
%----------------------------
\section{Non-Functional Requirements}
Cover scalability, availability, maintainability, reliability, latency, and security.

\subsection{Scalability}
Expected users, transactions, and data volume.

\subsection{Latency and Availability}
Define SLAs (e.g., P99, P50 latency targets).

\subsection{Maintainability}
Explain maintenance expectations and ownership.

\subsection{Security}
Describe security measures and required compliance levels.

\newpage
%----------------------------
\section{Testing and Observability}
Outline strategies for testing, monitoring, and alerting.

\subsection{Testing}
Explain testing approach: unit, integration, A/B, or stress tests.

\subsection{Metrics and Alarms}
List key performance metrics, dashboards, and alert configurations.

\newpage
%----------------------------
\section{Future Improvements}
List planned features or improvements for future releases.

\newpage
%----------------------------
\section{FAQs}
Include frequently asked questions and their answers.

question?
answer here

\newpage
%----------------------------
\appendix
\section{Appendix A: Subtitle (Optional)}
Include detailed technical analysis or supporting data.

\newpage
%----------------------------
\section{Glossary}
\begin{description}
  \item[ex1] example definition
\end{description}

\newpage
%----------------------------
\section{References}
\begin{itemize}
  \item Service 1: \url{https://example.com}
\end{itemize}

\end{document}
